% SP1 solution v3: adds counting argument (item 7); drops incorrect surjectivity claim from item 5.

I follow the hint.

For each \(s\in S\), choose a linear functional \(\phi_s:\mathbb{R}^2\to\mathbb{R}\) that is uniquely maximized at \(s\) on \(P=\operatorname{conv}(S)\). Since \(T\) is finite, choose \(t_s\in T\) with
\[
\phi_s(t_s)=\max_{t\in T}\phi_s(t),
\]
and define \(\ell_s:=s+t_s\).

\medskip
\textbf{1. Lonely points exist.}

Fix \(s\in S\). Suppose \(\ell_s = s'+t'\) for some \(s'\in S\), \(t'\in T\).
Applying \(\phi_s\):
\[
\phi_s(s)+\phi_s(t_s)=\phi_s(s')+\phi_s(t').
\]
By construction \(\phi_s(s')\le\phi_s(s)\) and \(\phi_s(t')\le\phi_s(t_s)\), so equality holds in both.
Since \(s\) is the unique maximizer of \(\phi_s\) on \(S\), we get \(s'=s\), hence \(t'=t_s\).
Therefore the only \(t\in T\) with \(\ell_s - t\in S\) is \(t=t_s\): the point \(\ell_s\) is lonely.

\medskip
\textbf{2. Distinctness.}

If \(\ell_{s_1}=\ell_{s_2}\), then by loneliness the covering translation is unique, so \(t_{s_1}=t_{s_2}\), hence \(s_1=s_2\).
Thus the \(n\) points \(\ell_s\) are pairwise distinct.

\medskip
\textbf{3. Non-vanishing.}

By loneliness, \(F(\ell_s)=\varepsilon(t_s)\neq 0\).

\medskip
\textbf{4. Lower bound.}

The \(n\) distinct points \(\ell_s\) all lie in \(\operatorname{supp}(F)\), so \(|\operatorname{supp}(F)|\ge n\).

\medskip
\textbf{5. Rigidity in the equality case.}

Assume \(|\operatorname{supp}(F)|=n\). Since \(\{\ell_s:s\in S\}\subseteq\operatorname{supp}(F)\) and both sets have size \(n\), we have
\[
\operatorname{supp}(F)=\{\ell_s:s\in S\}.
\]
In particular, for any \(p\in S+T\) that is not of the form \(\ell_s\), we have \(F(p)=0\).

\medskip
\textbf{6. Vertices of \(P+Q\) survive.}

Let \(v\) be a vertex of \(P+Q=\operatorname{conv}(S)+\operatorname{conv}(T)\).
Choose \(\psi:\mathbb{R}^2\to\mathbb{R}\) uniquely maximized at \(v\) on \(P+Q\).
Let \(s^*\in S\) and \(t^*\in T\) be the maximizers of \(\psi\) on \(S\) and \(T\) respectively.
Then \(\psi(s^*+t^*)=\max_{P+Q}\psi=\psi(v)\), so \(v=s^*+t^*\) by uniqueness.
The same uniqueness argument shows \(s^*\) is the unique maximizer of \(\psi\) on \(S\), and \(t^*\) is the unique maximizer on \(T\).
Hence \(v\) is lonely: the only pair \((s',t')\) with \(s'+t'=v\) is \((s^*,t^*)\).
Therefore \(F(v)=\varepsilon(t^*)\neq 0\), so every vertex of \(P+Q\) lies in \(\operatorname{supp}(F)\).

\medskip
\textbf{7. Sign-count identity (counting argument).}

Define \(\sigma(s):=F(\ell_s)=\varepsilon(t_s)\in\{+1,-1\}\) and let
\[
n_+:=|\{s\in S:\sigma(s)=+1\}|, \qquad n_-:=|\{s\in S:\sigma(s)=-1\}|.
\]
Also let \(k_1:=|\{t\in T:\varepsilon(t)=+1\}|\) and \(k_2:=|\{t\in T:\varepsilon(t)=-1\}|\).

We compute the total weight of \(F\) in two ways.

\emph{First way.} In the equality case \(\operatorname{supp}(F)=\{\ell_s:s\in S\}\), so
\[
\sum_{p\in\mathbb{R}^2}F(p)=\sum_{s\in S}F(\ell_s)=\sum_{s\in S}\sigma(s)=n_+-n_-.
\]

\emph{Second way.} By definition \(F(p)=\sum_{t\in T:\,p-t\in S}\varepsilon(t)\), so
\[
\sum_{p\in\mathbb{R}^2}F(p)=\sum_{t\in T}\varepsilon(t)\cdot|S|=n(k_1-k_2).
\]

Equating the two expressions:
\[
n_+-n_-=n(k_1-k_2).
\]
Since \(0\le n_\pm\le n\) and \(n_++n_-=n\), we have \(|n_+-n_-|\le n\), so \(n(k_1-k_2)\in\{-n,0,n\}\), giving:
\[
k_1-k_2\in\{-1,0,1\}.
\]
In particular:
\begin{itemize}
\item \textbf{Case~0} (\(k_1=k_2\)): \(n_+=n_-=n/2\). Both signs appear on \(S\), and \(n\) is even.
\item \textbf{Case~\(+n\)} (\(k_1=k_2+1\)): \(n_+=n\), \(n_-=0\). All lonely points have sign \(+1\).
\item \textbf{Case~\(-n\)} (\(k_2=k_1+1\)): \(n_-=n\), \(n_+=0\). All lonely points have sign \(-1\).
\end{itemize}
\end{proof}
